\documentclass[10pt,a4paper]{amsart}
\usepackage[margin=19mm]{geometry}
\usepackage{amsmath,amssymb,mathtools}
\usepackage[scr=rsfs,cal=euler]{mathalfa}
\usepackage{xfrac}
\usepackage[unicode,hidelinks,bookmarks=false]{hyperref}
\newcommand{\ie}{i.e.\ }
\newcommand{\cf}{cf.\ }
\newcommand{\resp}{respectively\ }
\newcommand{\Subsection}{\subsection}
\providecommand{\nobreakdash}{\nobreak\hbox{-}}
\setlength{\emergencystretch}{2em}
\multlinegap0pt
\usepackage{bm}
\usepackage{bbm}
\usepackage{dsfont}
\usepackage{xspace}
\usepackage{cancel}
\usepackage{mathtools}
\usepackage{amsthm}
\makeatletter
\@ifundefined{theorem}{\newtheorem{theorem}{Theorem}}{}
\@ifundefined{corollary}{\newtheorem{corollary}[theorem]{Corollary}}{}
\@ifundefined{definition}{\newtheorem{definition}[theorem]{Definition}}{}
\@ifundefined{proposition}{\newtheorem{proposition}[theorem]{Proposition}}{}
\makeatother
\makeatletter
\newcommand{\bZ}{{\mathbb{Z}}}
\newcommand{\mb}{\mathbf} 
\newcommand{\mc}{\mathcal} 
\expandafter\ifx\csname parts\endcsname\relax
\newenvironment{parts}[0]{%
  \begin{list}{}%
    {\setlength{\itemindent}{0pt}
     \setlength{\labelwidth}{1.5\parindent}
     \setlength{\labelsep}{.5\parindent}
     \setlength{\leftmargin}{2\parindent}
     \setlength{\itemsep}{0pt}
     }%
   }%
\fi
\newcommand{\torus}{{\mathbb{T}}}
\makeatother
\title[General Boyd-Lawton Theorems with Multivariable Limits]{General Boyd-Lawton Theorems with Multivariable Limits}
\author{Wayne Aitken, Kimberly Ayers,, Hanson Smith}
\date{Source: arXiv:2508.05910v1 (CC BY 4.0)}
\begin{document}
\maketitle

\noindent\emph{Source and licence.} General Boyd-Lawton Theorems with Multivariable Limits. Version arXiv:2508.05910v1 (CC BY 4.0), distributed under the Creative Commons Attribution 4.0 International licence, as recorded in the arXiv licence metadata for this exact version and shown on the abstract page https://arxiv.org/abs/2508.05910v1. Full rights notice: licenses/arxiv-2508.05910-CC-BY-4.0.txt.\par\smallskip

\noindent\textit{Editorial scope.} Counted unit: the Corollary whose source label is \texttt{cor\_mt} in the article (source0.tex lines 1601--1611, source label \texttt{cor\_mt}), delivered together with the article's own complete original proof of it (source0.tex lines 1538--1597, one \texttt{proof} environment of 60 lines). No mathematics is written, repaired or re-derived: statement and proof are the article's own text line for line. The proof is detached from the statement in the source: the article states the result at lines 1601--1611 and prints its proof at lines 1538--1597, and the proof environment's own header names the statement, so the association is the article's own. Required context retained uncounted: BL collection def (lines 504--524); infinite\_height\_lemma (lines 871--883). Every definition, display and label of those lines is retained unchanged. Omitted from this package and not redistributed: 1--503, 525--870, 884--1537, 1598--1600, 1612--1799. Nothing inside a retained excerpt is omitted or altered: every retained body is delivered exactly as the article's lines are written, so no omission receipt of any kind accompanies this package. Citations: the retained excerpts carry no citation command at all, so no bibliography entry of the article is transcribed and no reference is redistributed. Reference adaptation: none was needed and none was made; every reference inside the delivered unit resolves inside it, so no unresolved reference or citation is produced. Printed slips kept literal: no printed word, symbol, hypothesis or number of the counted statement or of its proof was altered. Layout and class: the article's own class is \texttt{amsart} with packages including amsmath, amsthm, amssymb, amscd, url, enumerate, mathtools, graphicx, geometry, afterpage, tikz, todonotes, xy, hyperref; the wrapper is the lane's pinned \texttt{amsart} editorial wrapper at 10pt on A4, which restates the theorem-like environments the retained text needs and adds the article's own macro definitions that the retained text needs (\texttt{\textbackslash bZ}, \texttt{\textbackslash mb}, \texttt{\textbackslash mc}, \texttt{\textbackslash torus}), with extra wrapper packages (bm, bbm, dsfont, xspace, cancel, mathtools). The delivered numbering is the unit's own. Assets: no retained span contains a figure environment, a graphics-inclusion command, a PDF-page inclusion, a table or a TikZ picture; no image file is copied into this package, the package declares no asset dependency and no image file is redistributed with it. Licence: version arXiv:2508.05910v1 of this article is distributed under the Creative Commons Attribution 4.0 International licence, as recorded in the arXiv licence metadata for this exact version and shown on the abstract page https://arxiv.org/abs/2508.05910v1. A full rights notice is delivered at licenses/arxiv-2508.05910-CC-BY-4.0.txt. Characters: every retained excerpt is pure ASCII, so the delivered engine, XeLaTeX with its default Latin Modern text font, reports no missing character.
\begin{definition} \label{BL collection def}
A \emph{Boyd--Lawton collection} 
 $\mc C = \bigcup_{n \ge 1} \mc C (\torus^n)$ is a collection of functions such that the following hold:
\begin{enumerate}
\item
If $g \in \mc C (\torus^n)$, then $g$ is defined and is complex-valued almost everywhere on $\torus^n$. Further, each $g \in \mc C(\torus^n)$
is integrable: it gives an element of the $L^1$-space~$L^1(\torus^n)$ where as usual $\torus^n$ is given the normalized Haar measure.
\item
For each $g\in \mc C(\torus^n)$ with $n\ge 2$, a classical Boyd--Lawton theorem holds:
$$
\int_{\torus^n}  g \, d\tau_n= \lim_{\mu(q_\mb r)\to \infty}\;  \int_\torus g \circ q_{\mb r} \, d \tau,
$$
where $q_{\mb r} \colon \torus \to \torus^n$ vary among the continuous homomorphisms $\torus \to\torus^n$ defined by lattice points~$\mb r \in \bZ^n$.
(Implicit here is the requirement that, for large enough $\mu\left(q_\mb r\right)$, the function $g \circ q_{\mb r}$ is integrable.)

\item
The class $\mc C$ is closed under composition by homomorphisms $q_{A}$
in the following sense: If $g$ is in $\mc C(\torus^n)$, then there is a bound $B$ such that if $q_A \colon \torus^m \to \torus^n$ has Boyd-height~$\mu (q_A) \ge B$, 
then the composition  $g\circ q_{A}$ is in $\mc C(\torus^m)$.
\end{enumerate}
\end{definition}

\begin{proposition} \label{infinite_height_lemma}
Let $q_A \colon \torus^m \to \torus^n$ be a continuous homomorphism where~$A \in M_{n, m}(\bZ)$.
Then the following are equivalent.
\begin{enumerate}
\item
The Boyd height $\mu\left( q_A \right) = \mu(A)$ is infinite.
\item
The matrix $A$ has rank $n$.
\item
The homomorphism $q_A \colon \torus^m \to \torus^n$ is surjective.
\end{enumerate}
In particular, a necessary condition for the height $\mu\left( q_A \right)$ to be infinite is $n \le m$.
\end{proposition}

\begin{corollary}  \label{cor_mt}
Let $\mc C$ be a Boyd--Lawton collection and let $g \in \mc C_n$.
Suppose $q_A \colon \torus^m \to \torus^n$ is a nonzero continuous homomorphism of infinite Boyd height. In other words, suppose 
that~$q_A \colon \torus^m \to \torus^n$ is surjective.
Then
$$
 \int_{\torus^{n}} g \, d \tau_{n}
 \;  = \; 
\int_{\torus^m} g \circ q_A \, d \tau_m \,.
$$
\end{corollary}

\begin{proof}[Proof for $n=1$.]
In this case we will prove something stronger: we will show that if~$g\in \mc C(\torus^1)$ and if 
$q_A \colon \torus^m \to \torus^1$ is a continuous homomorphism with $\mu\left(q_A\right) \ge M = 2$,
then
$$
 \int_{\torus^{1}} g \, d \tau
 \;  = \; 
\int_{\torus^{m}} g \circ q_{A} \ d \tau_{m}.
$$
Trivially the two integrals will be distance less than any $\varepsilon>0$.  Note: this equality of integrals is just a special case of Corollary~\ref{cor_mt} below.

Suppose $q_A \colon \torus^m \to \torus^1$ is a continuous homomorphism with~$\mu(A) = \mu\left(q_A\right) \ge 2$. Since~$\mu(A) \ne 1$, the row matrix $A$ is not the zero matrix, and so has rank $1=n$.
In particular,~$q_A \colon \torus^m \to \torus$ is surjective (Proposition~\ref{infinite_height_lemma}).

The main trick of this proof (for $n=1$) is to regard $g$ not as a function on $\torus^1$ as originally given, but as a function on~$\torus^2$ where $g$  simply ignores the second coordinate;
this sets us up to use the $n = 2$ case of the theorem which has already been established.
To make this precise, let $\tilde g$ be a function (almost everywhere) on $\torus^2= \torus \times \torus$
defined by the rule~$\tilde g(z_1, z_2) = g(z_1)$. 
In particular, using the Fubini-Tonelli theorem,
\begin{eqnarray*}
\int_{\torus^2} \tilde g \ d\tau_2 
&=& \int_{\torus^2} g (z_1) \ d\tau_2(z_1, z_2) \\
&=& \int_{\torus} \left( \int_{\torus} g(z_1) \ d\tau(z_2) \right) d\tau  (z_1) \\
&=& \int_{\torus}  g(z_1) \left( \int_{\torus}  \ d\tau \right) d\tau  (z_1)   = \int_{\torus} g \ d\tau.
\end{eqnarray*}
We also claim that $\tilde g$ is in $\mc C(\torus^2)$, where $\mc C$ is the  Boyd--Lawton collection given in the statement of the theorem.
To see this, note that  $\tilde g = g\circ q_{\mb t}$ on $\torus^2$ where $\mb t = [1, 0]$ and $q_{\mb t} \colon \torus^2 \to \torus$ is the projection map to the first coordinate. Since $q_{\mb t}$ is surjective, we have 
that $\mu \left(q_{\mb t}\right) = \infty$ (Proposition~\ref{infinite_height_lemma}), which forces $\tilde g= g\circ q_{\mb t}$ to be in $\mc C(\torus^2)$ by the definition of Boyd--Lawton collection (Definition~\ref{BL collection def}, part 3).


Next, we expand $q_A$ to a map $\torus^{m+1} \to \torus^2$ by having it map as the identity on the last coordinate. More precisely,
consider the homomorphism $q_{\tilde A} \colon \torus^m \times \torus \to \torus \times \torus$ given by
$$
q_{\tilde A} (\mb z, w) = \left( q_A(\mb z), w \right).
$$
Note that $q_{\tilde A}$ is continuous and,
since $q_A$ is surjective, $q_{\tilde A}$ must also be surjective.  This implies that~$\mu\left( q_{\tilde A} \right) = \infty$ (Proposition~\ref{infinite_height_lemma}).
Using the  Fubini-Tonelli theorem as before we have
\begin{eqnarray*}
\int_{\torus^{m+1}} \tilde g \circ q_{\tilde A} \ d\tau_{m+1} 
&=& \int_{\torus^m} \left(  \int_{\torus} \tilde g \circ q_{\tilde A} (\mb z, w) \ d\tau (w) \right) \ d\tau_{m} (\mb z) \\
&=& \int_{\torus^m} \left(  \int_{\torus} \tilde g  \left(q_A(\mb z), w\right) \ d\tau (w) \right) \ d\tau_{m} (\mb z) \\
&=& \int_{\torus^m} \left(  \int_{\torus} g  \left(q_A(\mb z)\right) \ d\tau (w) \right) \ d\tau_{m} (\mb z) \\
&=& \int_{\torus^m} g  \left(q_A(\mb z)\right)  \left( \int_{\torus}  \ d\tau \right) d\tau_m  (\mb z)   = \int_{\torus^m} g \circ q_A \ d\tau_m.
\end{eqnarray*}

Finally we note that we have established the theorem for $n>1$, so it can be applied to the function $\tilde g \in \mc C(\torus^2)$.
In fact, since   $\mu\left( q_{\tilde A} \right) = \infty$, we get the stronger result that
$$
 \int_{\torus^{2}} \tilde g \, d \tau_{n}
 \;  = \; 
\int_{\torus^{m+1}} \tilde g \circ q_{\tilde A} \ d \tau_{m+1},
$$
as highlighted in the corollary below.
By substitution, we have
\[
 \int_{\torus^{1}} g \, d \tau
 \;  = \; 
\int_{\torus^{m}} g \circ q_{A} \ d \tau_{m}. \qedhere\]
\end{proof} 

\end{document}
